\par\textbf{Procedimiento completo para comparar tu solución.} Continuidad obtiene velocidad de salida. Bernoulli horizontal calcula su presión. Solo entonces el balance de momento permite hallar la fuerza de pared sobre el fluido; la fuerza del fluido sobre el conducto tiene signo opuesto.

\begin{enumerate}

\item Conservar caudal volumétrico.

\[\fitmath{U2=\frac{0.12\cdot 22}{0.06}\approx 44\,\mathrm{m/s}}\]

\item Caudal másico de entrada.

\[\fitmath{\mathrm{mdot}=1.19\cdot 0.12\cdot 22\approx 3.1416\,\mathrm{kg/s}}\]

\item Despejar presión manométrica de salida.

\[\fitmath{p2=3200+\frac{1.19\cdot \left({22}^{2}-{U2}^{2}\right)}{2}}\]

\[\fitmath{p2=3200+\frac{1.19\cdot \left({22}^{2}-{(44)}^{2}\right)}{2}\approx 2336.06\,\mathrm{Pa}}\]

\item Fuerza neta de las dos caras de presión.

\[\fitmath{\mathrm{Fpressure}=3200\cdot 0.12-p2\cdot 0.06}\]

\[\fitmath{\mathrm{Fpressure}=3200\cdot 0.12-(2336.06)\cdot 0.06\approx 243.8364\,\mathrm{N}}\]

\item Restar las fuerzas de presión del cambio de flujo de momento.

\[\fitmath{\mathrm{Fwall}=\mathrm{mdot}\cdot \left(U2-22\right)-\mathrm{Fpressure}}\]

\[\fitmath{\mathrm{Fwall}=(3.1416)\cdot \left((44)-22\right)-(243.8364)\approx -174.7212\,\mathrm{N}}\]

\item Aplicar acción y reacción para la fuerza sobre el conducto.

\[\fitmath{\mathrm{Ffluid}=-\mathrm{Fwall}}\]

\[\fitmath{\mathrm{Ffluid}=-(-174.7212)\approx 174.7212\,\mathrm{N}}\]

\end{enumerate}
